\documentclass[exercice,a4]{thedocument}%
\startdatakeys
    \source{laclassecooperativedestef.com}
    \BO{2026}
    \cycle{3}
    \competencesDuSocle{RE4,CO2}
    \objectifsApprentissage{C3C1ab,C3C1aj,C3C1ae,C3C1af,C3C1ag}
    \motsCles{perpendicularite, parallelisme, quadrilatere}
\enddatakeys
\begin{document}%
\noindent\textbf{Programme de construction 27}\par\smallskip\noindent
\begin{itemize}
    \item Tracer une droite $(\mathspoint{d})$.
    \item Placer les points $\mathspoint{A}$, $\mathspoint{B}$,
    $\mathspoint{C}$ et $\mathspoint{D}$ sur $(\mathspoint{d})$ tels que
    $\mathspoint{AB}=1\text{ cm}$, $\mathspoint{BC}=5\text{ cm}$ et
    $\mathspoint{CD}=1\text{ cm}$.
    \item Tracer les droites perpendiculaires à $(\mathspoint{d})$ :
    \begin{itemize}
        \item[(a)] passant par $\mathspoint{A}$ ;
        \item[(b)] passant par $\mathspoint{B}$ ;
        \item[(c)] passant par $\mathspoint{C}$ ;
        \item[(e)] passant par $\mathspoint{D}$.
    \end{itemize}
    \item Tracer la droite $(\mathspoint{f})$ en dessous de
    $(\mathspoint{d})$, parallèle à $(\mathspoint{d})$ et distante de
    $4\text{ cm}$.
    \item La droite $(\mathspoint{f})$ coupe la droite (a) en
    $\mathspoint{E}$, la droite (b) en $\mathspoint{F}$, la droite (c) en
    $\mathspoint{G}$ et la droite (e) en $\mathspoint{H}$.
    \item Tracer le quadrilatère $\mathspoint{BDGE}$.
    \item Quel est le nom de cette figure ?
\end{itemize}
\end{document}%
